% answers.tex —— 2012 年·全国大纲卷·文科 · 参考答案与详细解析片段
% 由 tools/generate_projects.py 依来源总账自 2012/2012_数学真题_全国大纲文_试题.md + 2012/2012_数学真题_全国大纲文_解析正文.md + 2012/2012_数学真题_全国大纲文_答案要点.md 生成，勿手改（改源 MD 后重跑）


\begin{answer}
\gksolhead{1}{}

\textbf{题目}：

已知集合 $A=\{x\mid x\text{ \text{是平行四边形}}\}$，$B=\{x\mid x\text{ \text{是矩形}}\}$，$C=\{x\mid x\text{ \text{是正方形}}\}$，$D=\{x\mid x\text{ \text{是菱形}}\}$，则\mbox{（　　）}

（A）$A\subseteq B$　（B）$C\subseteq B$　（C）$D\subseteq C$　（D）$A\subseteq D$

\Needspace{4\baselineskip}
\textbf{解析}

正方形既是矩形又是菱形，故 $C\subseteq B$；其余三个包含关系均不总成立。因此选 B。

\gkanswer{答案：B}

\end{answer}
\begin{answer}
\gksolhead{2}{}

\textbf{题目}：

函数 $y=\sqrt{x+1}$（$x\geq-1$）的反函数为\mbox{（　　）}

（A）$y=x^2-1$（$x\geq0$）　（B）$y=x^2-1$（$x\geq1$）

（C）$y=x^2+1$（$x\geq0$）　（D）$y=x^2+1$（$x\geq1$）

\Needspace{4\baselineskip}
\textbf{解析}

\Needspace{5\baselineskip}
交换 $x,y$，得 $x=\sqrt{y+1}$。由原函数值域 $x\geq0$，平方后得反函数

\nopagebreak[4]
\[
y=x^2-1\quad(x\geq0).
\]

故选 A。

\gkanswer{答案：A}

\end{answer}
\begin{answer}
\gksolhead{3}{}

\textbf{题目}：

若函数 $f(x)=\sin\dfrac{x+\varphi}{3}$（$\varphi\in[0,2\pi]$）是偶函数，则 $\varphi=$\mbox{（　　）}

（A）$\dfrac\pi2$　（B）$\dfrac{2\pi}3$　（C）$\dfrac{3\pi}2$　（D）$\dfrac{5\pi}3$

\Needspace{4\baselineskip}
\textbf{解析}

\Needspace{5\baselineskip}
令 $u=(x+\varphi)/3$。要使 $\sin u$ 对 $x$ 为偶函数，须有相位

\nopagebreak[4]
\[
\frac{\varphi}{3}=\frac\pi2+k\pi.
\]

在 $\varphi\in[0,2\pi]$ 内只有 $\varphi=3\pi/2$，故选 C。

\gkanswer{答案：C}

\end{answer}
\begin{answer}
\gksolhead{4}{}

\textbf{题目}：

已知 $\alpha$ 为第二象限角，$\sin\alpha=\dfrac35$，则 $\sin2\alpha=$\mbox{（　　）}

（A）$-\dfrac{24}{25}$　（B）$-\dfrac{12}{25}$　（C）$\dfrac{12}{25}$　（D）$\dfrac{24}{25}$

\Needspace{4\baselineskip}
\textbf{解析}

\Needspace{5\baselineskip}
$\alpha$ 在第二象限，故

\nopagebreak[4]
\[
\cos\alpha=-\sqrt{1-\sin^2\alpha}=-\frac45.
\]

于是 $\sin2\alpha=2\sin\alpha\cos\alpha=-24/25$，选 A。

\gkanswer{答案：A}

\end{answer}
\begin{answer}
\gksolhead{5}{}

\textbf{题目}：

椭圆的中心在原点，焦距为 $4$，一条准线为 $x=-4$，则该椭圆的方程为\mbox{（　　）}

（A）$\dfrac{x^2}{16}+\dfrac{y^2}{12}=1$　（B）$\dfrac{x^2}{12}+\dfrac{y^2}{8}=1$

（C）$\dfrac{x^2}{8}+\dfrac{y^2}{4}=1$　（D）$\dfrac{x^2}{12}+\dfrac{y^2}{4}=1$

\Needspace{4\baselineskip}
\textbf{解析}

\Needspace{5\baselineskip}
焦距为 $4$，故 $c=2$。椭圆横轴的准线为 $x=\pm a^2/c$，由 $a^2/c=4$ 得 $a^2=8$，再由 $b^2=a^2-c^2$ 得 $b^2=4$。故方程为

\nopagebreak[4]
\[
\frac{x^2}{8}+\frac{y^2}{4}=1,
\]

选 C。

\gkanswer{答案：C}

\end{answer}
\begin{answer}
\gksolhead{6}{}

\textbf{题目}：

已知数列 $\{a_n\}$ 的前 $n$ 项和为 $S_n$，$a_1=1$，$S_n=2a_{n+1}$，则 $S_n=$\mbox{（　　）}

（A）$2^{n-1}$　（B）$\left(\dfrac32\right)^{n-1}$　（C）$\left(\dfrac23\right)^{n-1}$　（D）$\dfrac1{2^{n-1}}$

\Needspace{4\baselineskip}
\textbf{解析}

\Needspace{5\baselineskip}
由 $S_{n-1}=2a_n$，得 $a_n=S_{n-1}/2$，从而

\nopagebreak[4]
\[
S_n=S_{n-1}+a_n
=S_{n-1}+\frac{S_{n-1}}2
=\frac32S_{n-1}.
\]

\Needspace{5\baselineskip}
又 $S_1=a_1=1$，故

\nopagebreak[4]
\[
S_n=\left(\frac32\right)^{n-1},
\]

选 B。

\gkanswer{答案：B}

\end{answer}
\begin{answer}
\gksolhead{7}{}

\textbf{题目}：

$6$ 位选手依次演讲，其中选手甲不在第一个也不在最后一个演讲，则不同的演讲次序共有\mbox{（　　）}

（A）$240$ 种　（B）$360$ 种　（C）$480$ 种　（D）$720$ 种

\Needspace{4\baselineskip}
\textbf{解析}

\Needspace{5\baselineskip}
甲可排在中间 $4$ 个位置中的任一个，其余 $5$ 人任排，方法数为

\nopagebreak[4]
\[
4\times5!=480.
\]

故选 C。

\gkanswer{答案：C}

\end{answer}
\begin{answer}
\gksolhead{8}{}

\textbf{题目}：

已知正四棱柱 $ABCD-A_1B_1C_1D_1$ 中，$AB=2$，$CC_1=2\sqrt2$，$E$ 为 $CC_1$ 的中点，则直线 $AC_1$ 到平面 $BED$ 的距离为\mbox{（　　）}

（A）$2$　（B）$\sqrt3$　（C）$\sqrt2$　（D）$1$

\Needspace{4\baselineskip}
\textbf{解析}

\Needspace{5\baselineskip}
取 $A(0,0,0)$，$B(2,0,0)$，$C(2,2,0)$，$D(0,2,0)$，则

\nopagebreak[4]
\[
C_1=(2,2,2\sqrt2),\qquad E=(2,2,\sqrt2).
\]

\Needspace{5\baselineskip}
平面 $BED$ 的法向量可取 $(1,1,-\sqrt2)$，方程为

\nopagebreak[4]
\[
x+y-\sqrt2z-2=0.
\]

直线 $AC_1$ 的方向向量为 $(2,2,2\sqrt2)$，与该法向量垂直，故直线平行于平面；点 $A$ 到平面的距离为 $2/\sqrt4=1$，选 D。

\gkanswer{答案：D}

\end{answer}
\begin{answer}
\gksolhead{9}{}

\textbf{题目}：

$\triangle ABC$ 中，$AB$ 边的高为 $CD$，若 $\overrightarrow{CB}=\boldsymbol a$，$\overrightarrow{CA}=\boldsymbol b$，$\boldsymbol a\cdot\boldsymbol b=0$，$|\boldsymbol a|=1$，$|\boldsymbol b|=2$，则 $\overrightarrow{AD}=$\mbox{（　　）}

（A）$\dfrac13\boldsymbol a-\dfrac13\boldsymbol b$　（B）$\dfrac23\boldsymbol a-\dfrac23\boldsymbol b$　（C）$\dfrac35\boldsymbol a-\dfrac35\boldsymbol b$　（D）$\dfrac45\boldsymbol a-\dfrac45\boldsymbol b$

\Needspace{4\baselineskip}
\textbf{解析}

\Needspace{5\baselineskip}
以 $C$ 为原点，则 $A=\boldsymbol b$，$B=\boldsymbol a$。设 $D=\lambda\boldsymbol a+(1-\lambda)\boldsymbol b$。因 $CD\perp AB$，

\nopagebreak[4]
\[
D\cdot(\boldsymbol a-\boldsymbol b)=\lambda|\boldsymbol a|^2-(1-\lambda)|\boldsymbol b|^2=0.
\]

\Needspace{5\baselineskip}
即 $\lambda-4(1-\lambda)=0$，所以 $\lambda=4/5$。从而

\nopagebreak[4]
\[
\overrightarrow{AD}=D-\boldsymbol b=\frac45\boldsymbol a-\frac45\boldsymbol b,
\]

选 D。

\gkanswer{答案：D}

\end{answer}
\begin{answer}
\gksolhead{10}{}

\textbf{题目}：

已知 $F_1,F_2$ 为双曲线 $C:x^2-y^2=2$ 的左、右焦点，点 $P$ 在 $C$ 上，$|PF_1|=2|PF_2|$，则 $\cos\angle F_1PF_2=$\mbox{（　　）}

（A）$\dfrac14$　（B）$\dfrac35$　（C）$\dfrac34$　（D）$\dfrac45$

\Needspace{4\baselineskip}
\textbf{解析}

\Needspace{5\baselineskip}
双曲线的 $a=\sqrt2$，$c=2$，故 $F_1F_2=4$。设 $|PF_2|=d$，则由焦点距离之差为 $2a$ 得

\nopagebreak[4]
\[
2d-d=2\sqrt2,qquad |PF_2|=2\sqrt2,\quad |PF_1|=4\sqrt2.
\]

\Needspace{5\baselineskip}
在 $\triangle F_1PF_2$ 中用余弦定理：

\nopagebreak[4]
\[
\cos\angle F_1PF_2=
\frac{(4\sqrt2)^2+(2\sqrt2)^2-4^2}{2\cdot4\sqrt2\cdot2\sqrt2}=\frac34.
\]

故选 C。

\gkanswer{答案：C}

\end{answer}
\begin{answer}
\gksolhead{11}{}

\textbf{题目}：

已知 $x=\ln\pi$，$y=\log_5 2$，$z=\mathrm e^{-\frac12}$，则\mbox{（　　）}

（A）$x<y<z$　（B）$z<x<y$　（C）$z<y<x$　（D）$y<z<x$

\Needspace{4\baselineskip}
\textbf{解析}

\Needspace{5\baselineskip}
因为 $\ln\pi>\ln\mathrm e=1$，$\log_5 2<\log_5 5=1$；又 $2<\sqrt5$，所以 $\log_5 2<1/2$，而 $\mathrm e^{-1/2}>1/2$。并且 $\mathrm e^{-1/2}<1<\ln\pi$，故

\nopagebreak[4]
\[
y<z<x.
\]

选 D。

\gkanswer{答案：D}

\end{answer}
\begin{answer}
\gksolhead{12}{}

\textbf{题目}：

正方形 $ABCD$ 的边长为 $1$，点 $E$ 在边 $AB$ 上，点 $F$ 在边 $BC$ 上，$AE=BF=\dfrac13$。动点 $P$ 从 $E$ 出发沿直线向 $F$ 运动，每当碰到正方形的边时反弹，反弹时反射角等于入射角，当点 $P$ 第一次碰到 $E$ 时，$P$ 与正方形的边碰撞的次数为\mbox{（　　）}

（A）$8$　（B）$6$　（C）$4$　（D）$3$

\Needspace{4\baselineskip}
\textbf{解析}

把每次反射展开为平面镶嵌中的直线运动。以 $E=(1/3,0)$、$F=(1,1/3)$，展开直线的方向可取 $(2,1)$。第一次到达 $E$ 的等价像为 $(13/3,2)$，中间依次穿过 $4$ 条竖边和 $1$ 条横边，终点再次碰边也计一次。

因此总碰撞次数为 $4+1+1=6$，选 B。

\gkanswer{答案：B}

\end{answer}
\begin{answer}
\gksolhead{13}{}

\textbf{题目}：

$\left(x+\dfrac1{2x}\right)^8$ 的展开式中 $x^2$ 的系数为\underline{\hspace{2cm}}。

\Needspace{4\baselineskip}
\textbf{解析}

\Needspace{5\baselineskip}
二项展开的一般项为

\nopagebreak[4]
\[
\binom8k x^{8-k}\left(\frac1{2x}\right)^k
=\frac{\binom8k}{2^k}x^{8-2k}.
\]

令 $8-2k=2$，得 $k=3$，所求系数为 $\binom83/2^3=56/8=7$。

\gkanswer{答案：$7$}

\end{answer}
\begin{answer}
\gksolhead{14}{}

\textbf{题目}：

\Needspace{5\baselineskip}
若 $x,y$ 满足约束条件

\nopagebreak[4]
\[
\begin{cases}
x-y+1\geq0,\\
x+y-3\leq0,\\
x+3y-3\geq0,
\end{cases}
\]

则 $z=3x-y$ 的最小值为\underline{\hspace{2cm}}。

\Needspace{4\baselineskip}
\textbf{解析}

\Needspace{5\baselineskip}
可行域三个顶点分别为

\nopagebreak[4]
\[
(0,1),\qquad(1,2),\qquad(3,0).
\]

在这三点，$z=3x-y$ 分别为 $-1,1,9$，故最小值是 $-1$。

\gkanswer{答案：$-1$}

\end{answer}
\begin{answer}
\gksolhead{15}{}

\textbf{题目}：

当函数 $y=\sin x-\sqrt3\cos x$（$0\leq x<2\pi$）取得最大值时，$x=$\underline{\hspace{2cm}}。

\Needspace{4\baselineskip}
\textbf{解析}

\nopagebreak[4]
\[
\sin x-\sqrt3\cos x=2\sin\left(x-\frac\pi3\right).
\]

最大值为 $2$，当 $x-\pi/3=\pi/2$ 时取得，故 $x=5\pi/6$。

\gkanswer{答案：$\dfrac{5\pi}{6}$}

\end{answer}
\begin{answer}
\gksolhead{16}{}

\textbf{题目}：

已知正方体 $ABCD-A_1B_1C_1D_1$ 中，$E,F$ 分别为 $BB_1,CC_1$ 的中点，那么异面直线 $AE$ 与 $D_1F$ 所成角的余弦值为\underline{\hspace{2cm}}。

\Needspace{4\baselineskip}
\textbf{解析}

\Needspace{5\baselineskip}
设正方体棱长为 $1$，取

\nopagebreak[4]
\[
A(0,0,0),\quad E(1,0,1/2),\quad D_1(0,1,1),\quad F(1,1,1/2).
\]

\Needspace{5\baselineskip}
则两条异面直线的方向向量可取 $(1,0,1/2)$、$(1,0,-1/2)$，所以

\nopagebreak[4]
\[
\cos\theta=\frac{|1-1/4|}{\sqrt{5/4}\sqrt{5/4}}=\frac35.
\]

\gkanswer{答案：$\dfrac35$}

\end{answer}
\begin{answer}
\gksolhead{17}{}

\textbf{题目}：

$\triangle ABC$ 的内角 $A,B,C$ 成等差数列，其对边 $a,b,c$ 满足 $2b^2=3ac$，求 $A$。

\Needspace{4\baselineskip}
\textbf{解析}

\Needspace{5\baselineskip}
由 $A,B,C$ 成等差数列且 $A+B+C=\pi$，得 $B=\pi/3$。设 $A=\pi/3-d$，$C=\pi/3+d$。由正弦定理可设 $a=K\sin A$，$b=K\sin B$，$c=K\sin C$，原条件化为

\nopagebreak[4]
\[
2\sin^2\frac\pi3=3\sin\left(\frac\pi3-d\right)\sin\left(\frac\pi3+d\right).
\]

\Needspace{5\baselineskip}
利用 $\sin(u-d)\sin(u+d)=\sin^2u-\sin^2d$，得到

\nopagebreak[4]
\[
\sin^2d=\frac14.
\]

且 $|d|<\pi/3$，故 $d=\pm\pi/6$。因此 $A=\pi/6$ 或 $\pi/2$，即 $A=30^\circ$ 或 $90^\circ$。

\gkanswer{答案：$A=30^\circ$ 或 $A=90^\circ$。}

\end{answer}
\begin{answer}
\gksolhead{18}{}

\textbf{题目}：

已知数列 $\{a_n\}$ 中，$a_1=1$，前 $n$ 项和 $S_n=\dfrac{n+2}{3}a_n$。

（1）　求 $a_2,a_3$；

（2）　求 $\{a_n\}$ 的通项公式。

\Needspace{4\baselineskip}
\textbf{解析}

\gksubpart{(1)}

\Needspace{5\baselineskip}
分别令 $n=2,3$：

\nopagebreak[4]
\[
1+a_2=\frac43a_2,\qquad 1+a_2+a_3=\frac53a_3.
\]

故 $a_2=3$，$a_3=6$。

\gksubpart{(2)}

\Needspace{5\baselineskip}
当 $n\geq2$ 时，将

\nopagebreak[4]
\[
S_n=\frac{n+2}{3}a_n,\qquad S_{n-1}=\frac{n+1}{3}a_{n-1}
\]

\Needspace{5\baselineskip}
相减并用 $S_n-S_{n-1}=a_n$，得

\nopagebreak[4]
\[
(n-1)a_n=(n+1)a_{n-1}.
\]

\Needspace{5\baselineskip}
结合 $a_1=1$，连乘可得

\nopagebreak[4]
\[
a_n=\prod_{k=2}^n\frac{k+1}{k-1}=\frac{n(n+1)}2.
\]

\gkanswer{答案（1）：$a_2=3$，$a_3=6$。}

\gkanswer{答案（2）：$a_n=\dfrac{n(n+1)}2$（$n\in\mathbb N^+$）。}

\end{answer}
\begin{answer}
\gksolhead{19}{}

\textbf{题目}：

如图，四棱锥 $P-ABCD$ 中，底面 $ABCD$ 为菱形，$PA\perp$ 底面 $ABCD$，$AC=2\sqrt2$，$PA=2$，$E$ 是 $PC$ 上的一点，$PE=2EC$。
\nopagebreak[4]
\begin{center}% 2012 全国大纲卷文科第 19 题：四棱锥 P-ABCD。
% 实线表示原图可见棱，虚线表示被遮挡的底面对角线或高线。
\begin{tikzpicture}[scale=.92, line cap=round, line join=round]
  \coordinate (P) at (2.35,4.05);
  \coordinate (A) at (2.35,1.45);
  \coordinate (B) at (0.35,0.80);
  \coordinate (C) at (1.20,0.05);
  \coordinate (D) at (4.40,0.75);
  \coordinate (E) at ($(P)!0.6667!(C)$);

  % 底面菱形与侧棱
  \draw[thick] (B)--(C)--(D);
  \draw[dashed] (D)--(A)--(B) (B)--(D);
  \draw[thick] (P)--(B) (P)--(D) (P)--(C);
  \draw[dashed] (P)--(A);

  % 点 E 在 PC 上；E 与 B、D 的连线给出截面 BED。
  \draw[thick] (E)--(B) (E)--(D);

  \fill (P) circle (1.15pt) node[above] {$P$};
  \fill (A) circle (1.05pt) node[right] {$A$};
  \fill (B) circle (1.05pt) node[left] {$B$};
  \fill (C) circle (1.05pt) node[below] {$C$};
  \fill (D) circle (1.05pt) node[right] {$D$};
  \fill (E) circle (1.05pt) node[left] {$E$};
\end{tikzpicture}
\end{center}
\nopagebreak[4]
（1）　证明：$PC\perp$ 平面 $BED$；

（2）　设二面角 $A-PB-C$ 为 $90^\circ$，求 $PD$ 与平面 $PBC$ 所成角的大小。

\Needspace{4\baselineskip}
\textbf{解析}

\Needspace{5\baselineskip}
取底面为 $z=0$，令

\nopagebreak[4]
\[
A(0,0,0),\quad P(0,0,2),\quad C(2\sqrt2,0,0).
\]

\Needspace{5\baselineskip}
设 $B=(\sqrt2,-t,0)$，$D=(\sqrt2,t,0)$；由 $PE:EC=2:1$，

\nopagebreak[4]
\[
E=\left(\frac{4\sqrt2}{3},0,\frac23\right).
\]

\gksubpart{(1)}

\Needspace{5\baselineskip}
有

\nopagebreak[4]
\[
\overrightarrow{PC}=(2\sqrt2,0,-2),
\]

\Needspace{5\baselineskip}
并且

\nopagebreak[4]
\[
\overrightarrow{PC}\cdot\overrightarrow{EB}
=\overrightarrow{PC}\cdot\overrightarrow{ED}=0.
\]

直线 $EB,ED$ 相交且都在平面 $BED$ 内，故 $PC\perp$ 平面 $BED$。

\gksubpart{(2)}

\Needspace{5\baselineskip}
二面角 $A-PB-C$ 为直角，等价于平面 $APB$ 与 $CPB$ 的法向量垂直。由坐标叉积得到

\nopagebreak[4]
\[
4t^2-8=0,qquad t=\sqrt2.
\]

\Needspace{5\baselineskip}
此时平面 $PBC$ 的一个法向量可取 $(-2\sqrt2,2\sqrt2,-4)$，且

\nopagebreak[4]
\[
\overrightarrow{PD}=(\sqrt2,\sqrt2,-2).
\]

\Needspace{5\baselineskip}
于是线面角 $\theta$ 满足

\nopagebreak[4]
\[
\sin\theta=
\frac{|\overrightarrow{PD}\cdot(-2\sqrt2,2\sqrt2,-4)|}
{|PD|\sqrt{(-2\sqrt2)^2+(2\sqrt2)^2+(-4)^2}}
=\frac12.
\]

所以 $\theta=30^\circ$。

\gkanswer{答案（1）：见解析版}

\gkanswer{答案（2）：$PD$ 与平面 $PBC$ 所成角为 $30^\circ$。}

\end{answer}
\begin{answer}
\gksolhead{20}{}

\textbf{题目}：

乒乓球比赛规则规定：一局比赛，双方比分在 $10$ 平前，一方连续发球 $2$ 次后，对方再连续发球 $2$ 次，依次轮换。每次发球，胜方得 $1$ 分，负方得 $0$ 分。设在甲、乙的比赛中，每次发球，发球方得 $1$ 分的概率为 $0.6$，各次发球的胜负结果相互独立。甲、乙的一局比赛中，甲先发球。

（1）　求开始第 $4$ 次发球时，甲、乙的比分为 $1$ 比 $2$ 的概率；

（2）　求开始第 $5$ 次发球时，甲得分领先的概率。

\Needspace{4\baselineskip}
\textbf{解析}

\gksubpart{(1)}

\Needspace{5\baselineskip}
开始第 $4$ 次发球前，前三次发球依次由甲、甲、乙进行。比分为 $1:2$ 要求：甲的两次发球中恰一次得分，而乙第 $3$ 次发球得分。因此概率为

\nopagebreak[4]
\[
2\times0.6\times0.4\times0.6=0.288=\frac{36}{125}.
\]

\gksubpart{(2)}

\Needspace{5\baselineskip}
前两次甲发球所得的甲方分数记为 $X$，后两次乙发球所得的乙方分数记为 $Y$。则 $X,Y$ 独立同分布，且均服从 $B(2,0.6)$。四次后甲、乙得分之差为

\nopagebreak[4]
\[
[X+(2-Y)]-[(2-X)+Y]=2(X-Y).
\]

\Needspace{5\baselineskip}
故甲领先等价于 $X>Y$。由对称性，

\nopagebreak[4]
\[
P(X>Y)=\frac{1-P(X=Y)}2.
\]

\Needspace{5\baselineskip}
其中

\nopagebreak[4]
\[
P(X=Y)=0.16^2+0.48^2+0.36^2=0.3856,
\]

故所求概率为 $(1-0.3856)/2=0.3072=192/625$。

\gkanswer{答案（1）：$\dfrac{36}{125}$。}

\gkanswer{答案（2）：$\dfrac{192}{625}$。}

\end{answer}
\begin{answer}
\gksolhead{21}{}

\textbf{题目}：

已知函数 $f(x)=\dfrac13x^3+x^2+ax$。

（1）　讨论 $f(x)$ 的单调性；

（2）　设 $f(x)$ 有两个极值点 $x_1,x_2$，若过两点 $(x_1,f(x_1))$，$(x_2,f(x_2))$ 的直线 $l$ 与 $x$ 轴的交点在曲线 $f(x)$ 上，求 $a$ 的值。

\Needspace{4\baselineskip}
\textbf{解析}

\gksubpart{(1)}

\nopagebreak[4]
\[
f'(x)=x^2+2x+a=(x+1)^2+a.
\]

当 $a\geq0$ 时，$f'(x)\geq0$，故 $f$ 在 $\mathbb R$ 上单调递增。若 $a<0$，令 $s=\sqrt{-a}$，则导数零点为 $-1-s,-1+s$，其符号依次为正、负、正，故单调区间如答案要点所列。

\gksubpart{(2)}

\Needspace{5\baselineskip}
两极值点要求 $a<0$。令 $s=\sqrt{-a}$，并令 $u=x+1$，则

\nopagebreak[4]
\[
f(x)=\frac13u^3-(1+s^2)u+\frac23+s^2.
\]

\Needspace{5\baselineskip}
极值点对应 $u=\pm s$。两点连线的斜率为 $-1-2s^2/3$，它与 $x$ 轴的交点横坐标为

\nopagebreak[4]
\[
x_0=\frac{s^2-1}{2s^2+3}.
\]

题设要求 $f(x_0)=0$。当 $s^2>0$ 时，代入后除 $s^2=1$ 所给的 $x_0=0$ 外，余因子恒为负，不能为零；故 $s^2=1$，从而 $a=-1$。

\gkanswer{答案（1）：$a\geq0$ 时 $f$ 在 $\mathbb R$ 上单调递增；$a<0$ 时，令 $s=\sqrt{-a}$，则 $f$ 在 $(-\infty,-1-s)$、$(-1+s,+\infty)$ 上递增，在 $(-1-s,-1+s)$ 上递减。}

\gkanswer{答案（2）：$a=-1$。}

\end{answer}
\begin{answer}
\gksolhead{22}{}

\textbf{题目}：

已知抛物线 $C:y=(x+1)^2$ 与圆 $M:(x-1)^2+\left(y-\dfrac12\right)^2=r^2$（$r>0$）有一个公共点 $A$，且在 $A$ 处两曲线的切线为同一直线 $l$。

（1）　求 $r$；

（2）　设 $m,n$ 是异于 $l$ 且与 $C$ 及 $M$ 都相切的两条直线，$m,n$ 的交点为 $D$，求 $D$ 到 $l$ 的距离。

\Needspace{4\baselineskip}
\textbf{解析}

\gksubpart{(1)}

\Needspace{5\baselineskip}
设切点为 $A(t-1,t^2)$。抛物线在 $A$ 的切线方向向量为 $(1,2t)$。圆心为 $M(1,1/2)$，相切时 $\overrightarrow{MA}$ 应垂直该方向，故

\nopagebreak[4]
\[
(t-2)+2t\left(t^2-\frac12\right)=0
\quad\Longrightarrow\quad 2t^3-2=0.
\]

\Needspace{5\baselineskip}
得 $t=1$，即 $A=(0,1)$。所以

\nopagebreak[4]
\[
r=|MA|=\sqrt{1+\frac14}=\frac{\sqrt5}{2},
\]

且公共切线为 $l:y=2x+1$。

\gksubpart{(2)}

\Needspace{5\baselineskip}
抛物线参数为 $t$ 处的切线为

\nopagebreak[4]
\[
y=2tx+2t-t^2.
\]

\Needspace{5\baselineskip}
它与圆相切的条件是圆心到该直线的距离为 $\sqrt5/2$。平方并整理得

\nopagebreak[4]
\[
(t-1)^2(t^2-6t-1)=0.
\]

\Needspace{5\baselineskip}
$t=1$ 对应 $l$；另两条公切线对应

\nopagebreak[4]
\[
t_1=3-\sqrt{10},\qquad t_2=3+\sqrt{10}.
\]

\Needspace{5\baselineskip}
将这两条切线联立，得到 $D=(2,-1)$。因此

\nopagebreak[4]
\[
d(D,l)=\frac{|2\cdot2-(-1)+1|}{\sqrt{2^2+(-1)^2}}
=\frac6{\sqrt5}=\frac{6\sqrt5}{5}.
\]

\gkanswer{答案（1）：$r=\dfrac{\sqrt5}{2}$。}

\gkanswer{答案（2）：$\dfrac{6\sqrt5}{5}$。}

\end{answer}